% GENERATED FILE. Edit canonical lessons, then run npm run generate:books.
% canonical-source-hash: fnv1a64:ac47a6c146759377
% canonical-lessons: MW-C231-bako, MW-C231-sidho, MW-C231-ubho, MW-C231-gol, MW-C231-dhilo

\chapter{Bent and Straight}
\label{ch:mw-bent-and-straight}
\hlchaptermodality{231}

\begin{chapteropening}
\textbf{By the end of this chapter:} \emph{I can say whether a path is crooked or straight, and how a thing stands.}

Complete the last lesson of chapter 231: I can say whether a path is crooked or straight, and how a thing stands.
\end{chapteropening}

\section[\texorpdfstring{\mw{बांको} (bā̃ko)}{बांको (bā̃ko)}]{\mw{बांको} (bā̃ko) — crooked; bent}

\label{lesson:MW-C231-bako}



\begin{quote}
Before the new one: say the Marwadi for swallow-wort, then the Marwadi for a thorn.
\end{quote}

\subsection*{You'll want to know: \mw{बांको}}
\textbf{\mw{बांको}} — \emph{bā̃ko} — \textquotedblleft{}crooked; bent\textquotedblright{}.

\begin{quote}
\mw{डूंगर} \mw{माथै} \mw{गेलो} \mw{बांको} \mw{है।}
\end{quote}

\begin{quote}
\emph{ḍū̃gar māthai gelo bā̃ko hai.}
\end{quote}

\begin{quote}
On the hill the path is crooked.
\end{quote}

\begin{cousinweb}
\textbf{\mw{बांको}} (\emph{bā̃ko}) is inherited from Sanskrit \emph{vaṅka}, \textquotedblleft{}bent, crooked\textquotedblright{}; CDIAL lists Marwari \emph{bā̃ko} (CDIAL 11191).
\end{cousinweb}

\begin{grammarlens}[title={What you've built}]
One more word for this chapter.
\end{grammarlens}

\subsection*{Your turn}
\begin{itemize}
\raggedright
  \item \emph{Say it:} \emph{bā̃ko}
  \item \emph{Say it:} \emph{bā̃ko}, once more
  \item \emph{Say it:} \emph{ḍū̃gar māthai gelo bā̃ko hai.} — On the hill the path is crooked.
  \item \emph{Recall:} say the Marwadi for swallow-wort, then the Marwadi for a thorn, then say \emph{bā̃ko} again
\end{itemize}

\begin{tcolorbox}[breakable,colback=teal!4,colframe=teal!35!black,title={Before you move on}]
What does \mw{बांको} mean? (\textquotedblleft{}Crooked; bent\textquotedblright{}.) Say it once more.
\end{tcolorbox}
\section[\texorpdfstring{\mw{सीधो} (sīdho)}{सीधो (sīdho)}]{\mw{सीधो} (sīdho) — straight}

\label{lesson:MW-C231-sidho}



\begin{quote}
Before the new one: say the Marwadi for a thorn, then the Marwadi for crooked.
\end{quote}

\subsection*{You'll want to know: \mw{सीधो}}
\textbf{\mw{सीधो}} — \emph{sīdho} — \textquotedblleft{}straight\textquotedblright{}.

\begin{quote}
\mw{सड़क} \mw{सीधी} \mw{है।}
\end{quote}

\begin{quote}
\emph{saṛak sīdhī hai.}
\end{quote}

\begin{quote}
The road is straight.
\end{quote}

\begin{cousinweb}
\textbf{\mw{सीधो}} (\emph{sīdho}) is inherited from Sanskrit \emph{siddha}, \textquotedblleft{}accomplished\textquotedblright{}; CDIAL lists Hindi \emph{sīdhā}, \textquotedblleft{}straight\textquotedblright{}, and Gujarati \emph{sīdhũ}, \textquotedblleft{}straightforward\textquotedblright{}. CDIAL gives no Marwari form in this entry (CDIAL 13401).
\end{cousinweb}

\begin{grammarlens}[title={What you've built}]
One more word for this chapter.
\end{grammarlens}

\subsection*{Your turn}
\begin{itemize}
\raggedright
  \item \emph{Say it:} \emph{sīdho}
  \item \emph{Say it:} \emph{sīdho}, once more
  \item \emph{Say it:} \emph{saṛak sīdhī hai.} — The road is straight.
  \item \emph{Recall:} say the Marwadi for a thorn, then the Marwadi for crooked, then say \emph{sīdho} again
\end{itemize}

\begin{tcolorbox}[breakable,colback=teal!4,colframe=teal!35!black,title={Before you move on}]
What does \mw{सीधो} mean? (\textquotedblleft{}Straight\textquotedblright{}.) Say it once more.
\end{tcolorbox}
\section[\texorpdfstring{\mw{ऊभो} (ūbho)}{ऊभो (ūbho)}]{\mw{ऊभो} (ūbho) — standing; upright}

\label{lesson:MW-C231-ubho}



\begin{quote}
Before the new one: say the Marwadi for crooked, then the Marwadi for straight.
\end{quote}

\subsection*{You'll want to know: \mw{ऊभो}}
\textbf{\mw{ऊभो}} — \emph{ūbho} — \textquotedblleft{}standing; upright\textquotedblright{}.

\begin{quote}
\mw{ऊंट} \mw{ऊभो} \mw{है।}
\end{quote}

\begin{quote}
\emph{ū̃ṭ ūbho hai.}
\end{quote}

\begin{quote}
The camel is standing.
\end{quote}

\begin{cousinweb}
\textbf{\mw{ऊभो}} (\emph{ūbho}) is inherited from Sanskrit \emph{ūrdhva}, \textquotedblleft{}erect\textquotedblright{}; CDIAL lists Old Marwari \emph{ūbhau}, \textquotedblleft{}upright\textquotedblright{} (CDIAL 2426).
\end{cousinweb}

\begin{grammarlens}[title={What you've built}]
One more word for this chapter.
\end{grammarlens}

\subsection*{Your turn}
\begin{itemize}
\raggedright
  \item \emph{Say it:} \emph{ūbho}
  \item \emph{Say it:} \emph{ūbho}, once more
  \item \emph{Say it:} \emph{ū̃ṭ ūbho hai.} — The camel is standing.
  \item \emph{Recall:} say the Marwadi for crooked, then the Marwadi for straight, then say \emph{ūbho} again
\end{itemize}

\begin{tcolorbox}[breakable,colback=teal!4,colframe=teal!35!black,title={Before you move on}]
What does \mw{ऊभो} mean? (\textquotedblleft{}Standing; upright\textquotedblright{}.) Say it once more.
\end{tcolorbox}
\section[\texorpdfstring{\mw{गोळ} (goḷ)}{गोळ (goḷ)}]{\mw{गोळ} (goḷ) — round}

\label{lesson:MW-C231-gol}



\begin{quote}
Before the new one: say the Marwadi for straight, then the Marwadi for standing.
\end{quote}

\subsection*{You'll want to know: \mw{गोळ}}
\textbf{\mw{गोळ}} — \emph{goḷ} — \textquotedblleft{}round\textquotedblright{}.

\begin{quote}
\mw{आज} \mw{चांद} \mw{गोळ} \mw{है।}
\end{quote}

\begin{quote}
\emph{āj chā̃d goḷ hai.}
\end{quote}

\begin{quote}
Today the moon is round.
\end{quote}

\begin{cousinweb}
\textbf{\mw{गोळ}} (\emph{goḷ}) is inherited from Sanskrit \emph{gola}, \textquotedblleft{}a ball\textquotedblright{}; CDIAL lists Hindi \emph{gol}, \textquotedblleft{}anything round\textquotedblright{}, and Gujarati \emph{goḷ}, \textquotedblleft{}round\textquotedblright{}. CDIAL gives no Marwari form in this entry (CDIAL 4321).
\end{cousinweb}

\begin{grammarlens}[title={What you've built}]
One more word for this chapter.
\end{grammarlens}

\subsection*{Your turn}
\begin{itemize}
\raggedright
  \item \emph{Say it:} \emph{goḷ}
  \item \emph{Say it:} \emph{goḷ}, once more
  \item \emph{Say it:} \emph{āj chā̃d goḷ hai.} — Today the moon is round.
  \item \emph{Recall:} say the Marwadi for straight, then the Marwadi for standing, then say \emph{goḷ} again
\end{itemize}

\begin{tcolorbox}[breakable,colback=teal!4,colframe=teal!35!black,title={Before you move on}]
What does \mw{गोळ} mean? (\textquotedblleft{}Round\textquotedblright{}.) Say it once more.
\end{tcolorbox}
\section[\texorpdfstring{\mw{ढीलो} (ḍhīlo)}{ढीलो (ḍhīlo)}]{\mw{ढीलो} (ḍhīlo) — loose; slack}

\label{lesson:MW-C231-dhilo}



\begin{quote}
Before the new one: say the Marwadi for standing, then the Marwadi for round.
\end{quote}

\subsection*{You'll want to know: \mw{ढीलो}}
\textbf{\mw{ढीलो}} — \emph{ḍhīlo} — \textquotedblleft{}loose; slack\textquotedblright{}.

\begin{quote}
\mw{रस्सी} \mw{ढीली} \mw{है।}
\end{quote}

\begin{quote}
\emph{rassī ḍhīlī hai.}
\end{quote}

\begin{quote}
The rope is loose.
\end{quote}

\begin{cousinweb}
\textbf{\mw{ढीलो}} (\emph{ḍhīlo}) is inherited from a reconstructed \emph{*ḍhila}, \textquotedblleft{}loose, slack\textquotedblright{} (CDIAL 5590, which lists Marwari \emph{ḍhīlo}).
\end{cousinweb}

\begin{grammarlens}[title={What you've built}]
One more word for this chapter.
\end{grammarlens}

\subsection*{Your turn}
\begin{itemize}
\raggedright
  \item \emph{Say it:} \emph{ḍhīlo}
  \item \emph{Say it:} \emph{ḍhīlo}, once more
  \item \emph{Say it:} \emph{rassī ḍhīlī hai.} — The rope is loose.
  \item \emph{Recall:} say the Marwadi for standing, then the Marwadi for round, then say \emph{ḍhīlo} again
\end{itemize}

\begin{tcolorbox}[breakable,colback=teal!4,colframe=teal!35!black,title={Before you move on}]
What does \mw{ढीलो} mean? (\textquotedblleft{}Loose; slack\textquotedblright{}.) Say it once more.
\end{tcolorbox}
